\subsection{Why the whole sequence must be checked}
Let $\mathcal C_i$ contain the permitted lane projections for observation $i$, each with its lane, arc position and projection distance. A pointwise check asks only whether $\mathcal C_i$ is nonempty. A sequence check seeks $b_i\in\mathcal C_i$ for all $i$ with compatible transitions from $b_i$ to $b_{i+1}$. The two tests can give very different answers even when the input coordinates look plausible on a map.

Consider two trajectories near the same intersection. One passes through the junction on a connected internal lane. The second crosses from a bridge to the street below. Both might stay close to a road at all times. Only the first admits a connected sequence of states through the crossing. A similar issue arises with direction: traveling along a one-way lane backwards can have zero geometric error at every point. Direction and connectivity therefore enter the acceptance decision before it is meaningful to compare dynamic realism on a road subset.

The implementation maintains partial assignments over successive candidate sets. A transition can continue on one directed lane or follow graph connections, subject to a bound on plausible route distance for that step. The audit also allows a permissive approximation to an adjacent lane change on the same normal edge. The stricter reconstruction used for diffusion targets excludes lateral lane switches because their arc progress is ambiguous in the present representation. Eligibility and audit thus use related but different rules.

The relation between reconstruction and decoding can be checked on the observed training data. Projecting the strict training paths back to coordinates gave mean position discrepancy of 0.277 metres over the 425 reconstructed tracks. The analogous validation reconstruction error was 0.267 metres over 57 tracks. These are reconstruction diagnostics using the observed route. They measure the geometric approximation available to the representation on its selected population. Supplying such a reconstructed future path for held-out generation would answer an easier problem, controlled replay of known geometry.
